\documentclass{bxjsreport}
\usepackage[utf8]{inputenc}
\usepackage[T1]{fontenc}
\usepackage{geometry}
\usepackage{graphicx}
\usepackage{float}
\usepackage{hyperref}
\usepackage[english,indonesian]{babel}

\title{Optimasi Pemrosesan Template LaTeX Menggunakan Pendekatan Berbasis AI}
\author{Academic Support Agent}
\date{\today}

\begin{document}
\maketitle

\begin{abstract}

\end{abstract}

\chapter{PENDAHULUAN}


\chapter{TINJAUAN PUSTAKA}


% Original content (commented out):
% \documentclass[12pt, xelatex,ja=standard,jafont=noto]{bxjsreport}

\usepackage{playful_masterthesis}
\usepackage{natbib}
\usepackage{graphicx}
\usepackage{amsmath}
\usepackage{hyperref}
\hypersetup{
    colorlinks=false
}
\CJKspace

%----
\year{2021} %修了年月

\affiliation{立命館大学大学院 情報理工学研究科} % 大学名
\school{博士前期課程 情報理工学専攻\\人間情報コース} % 研究科と専攻

\supervisor{松村 耕平 准教授} % 指導教官
\setcounter{numadvisor}{0} % 審査委員の数（主査含む）
\advisor{野間 春生 教授} % 主査
\advisora{野間 春生1 教授} % 副査1
\advisorb{野間 春生2 教授} % 副査2
\advisorc{野間 春生3 教授} % 副査3
\advisord{野間 春生4 教授} % 副査4

\title{非ユークリッド空間における\\球の充填と被覆の問題について} % 論文タイトル
\author{松村 耕平} % 氏名
\authorid{66121500710} % 学籍番号

%----本文はここから---

\begin{document}
\maketitle
\begin{abstract}
本論文では ...
本論文では ...
本論文では ...
本論文では ...本論文では ...本論文では ...本論文では ...本論文では ...本論文では ...本論文では ...本論文では ...本論文では ...本論文では ...本論文では ...本論文では ...本論文では ...本論文では ...本論文では ...本論文では ...本論文では ...本論文では ...本論文では ...本論文では ...本論文では ...本論文では ...本論文では ...本論文では ...本論文では ...本論文では ...本論文では ...本論文では ...本論文では ...本論文では ...本論文では ...本論文では ...本論文では ...本論文では ...本論文では ...本論文では ...本論文では ...本論文では ...本論文では ...本論文では ...本論文では ...本論文では ...本論文では ...本論文では ...本論文では ...本論文では ...本論文では ...本論文では ...本論文では ...本論文では ...本論文では ...本論文では ...本論文では ...本論文では ...本論文では ...本論文では ...本論文では ...本論文では ...本論文では ...本論文では ...本論文では ...本論文では ...本論文では ...本論文では ...本論文では ...本論文では ...本論文では ...本論文では ...本論文では ...本論文では ...本論文では ...本論文では ...本論文では ...本論文では ...本論文では ...本論文では ...本論文では ...本論文では ...本論文では ...本論文では ...本論文では ...本論文では ...本論文では ...本論文では ...
\end{abstract}

%\begin{eabstract}
%This paper describes ...
%\end{eabstract}

\tableofcontents


\chapter{はじめに}
ぞうの卵はおいしいぞう。ぞうの卵はおいしいぞう。ぞうの卵はおいしいぞう。ぞうの卵はおいしいぞう。
\section{セクション1}
ぞうの卵はおいしいぞう。ぞうの卵はおいしいぞう。ぞうの卵はおいしいぞう。ぞうの卵はおいしいぞう。
\subsection{サブセクション1}
ぞうの卵はおいしいぞう。ぞうの卵はおいしいぞう。ぞうの卵はおいしいぞう。ぞうの卵はおいしいぞう。
\subsubsection{サブサブセクション1}
ぞうの卵はおいしいぞう。ぞうの卵はおいしいぞう。ぞうの卵はおいしいぞう。ぞうの卵はおいしいぞう。
\paragraph{パラグラフ}
図\ref{fig:lena}に示したように、ぞうの卵はおいしいぞう。ぞうの卵はおいしいぞう。ぞうの卵はおいしいぞう。ぞうの卵はおいしいぞう。

\begin{figure}[htbp]
\centering
\includegraphics[scale=0.2]{images/lena.jpg}
\caption{A standard test image widely used in the field of image processing.}
\label{fig:lena}
\end{figure}

\chapter{関連研究}
ぞうの卵はおいしいぞう。ぞうの卵はおいしいぞう。ぞうの卵はおいしいぞう。ぞうの卵はおいしいぞう\citep{6926269,Matsumura:2013:CAR:2468356.2468600, Matsumura:2014:CFS:2582051.2582068,Matsumura:2014:YTW:2667317.2667417,matsumura2011role,matsumura2012office}。
\chapter{アプローチ}
ぞうの卵はおいしいぞう。ぞうの卵はおいしいぞう。ぞうの卵はおいしいぞう。ぞうの卵はおいしいぞう。
\chapter{実験}
ぞうの卵はおいしいぞう。ぞうの卵はおいしいぞう。ぞうの卵はおいしいぞう。ぞうの卵はおいしいぞう。
\chapter{議論}
ぞうの卵はおいしいぞう。ぞうの卵はおいしいぞう。ぞうの卵はおいしいぞう。ぞうの卵はおいしいぞう。
\chapter{結論}
ぞうの卵はおいしいぞう。ぞうの卵はおいしいぞう。ぞうの卵はおいしいぞう。ぞうの卵はおいしいぞう。


% 参考文献はbibtexで
\bibliographystyle{plain}
\bibliography{references}

\chapter*{謝辞}
お父さん、お母さん、今まで育ててくれてありがとう。ぞうの卵はおいしいぞう。ぞうの卵はおいしいぞう。ぞうの卵はおいしいぞう。ぞうの卵はおいしいぞう。
\end{document}


\end{document}